\section{S.M.A.R.T. goals}
\label{ResearchSmartGoals}
The Goals of this project are now defined here as “\gls{SmartGoals}”, as shown in the following list. They are all \textbf{S}pecific: the goal is clearly formulated; \textbf{M}easurable: there are measurable criteria defining when the goal is reached; \textbf{A}chievable: they are realistic given the available resources; \textbf{R}elevant: the goal is directly linked to the project’s success; and \textbf{T}ime-bound: there is a defined timeframe for achievement.


\begin{itemize}
\item \textbf{1) Setup of an improved data quality report}
    
\textbf{WHY:} This will help the team to run the daily checks faster

By the end of the implementation phase, the existing data quality report is to be updated to use the new infrastructure. This includes new reporting data pipelines and the new structure using star shema. Also, Power BI is to be used.

Success is measured by user feedback which has to be retrieved by \gls{ASTV} team. The feedback is to be implemented and the report must be on the productive systems.

..................
Todo
Rewrite
GPT Stuff
...............



\end{itemize}



SMART Goal 2 – Evaluation of a future DQM platform

WHY: To identify a suitable technological foundation for future enterprise-wide data quality management.

By the end of the evaluation phase, available DQM solutions shall be compared against the technical and organisational requirements identified within Swiss Post.

Success is measured by a documented comparison of the evaluated tools, including their strengths, limitations, and an evidence-based recommendation for future implementation.

SMART Goal 3 – Validation of the selected monitoring technology

WHY: To verify that the selected technology can support the implementation of automated data quality monitoring in future work.

By the end of the implementation phase, a proof of concept shall demonstrate that the selected monitoring solution can retrieve representative data, process the returned information, and provide a suitable basis for automated DQM checks.

Success is measured by the successful execution of the proof of concept, including documented test results and an assessment of its suitability for the subsequent Master Thesis.